\subsubsection{Finite normal covers of the base and iterated fibrations}

For a normal connected proper scheme $B$ over an algebraically
closed field, consider the following property:
\[
\begin{gathered}
\mathcal H(B):\quad\text{for every finite dominant morphism
$B'\to B$ with $B'$ normal and connected,}\\
\text{and every proper closed subset $T\subset B'$, the group
$Q(B',T)$ is topologically finitely generated.}
\end{gathered}
\]
The finite morphisms in this definition need not be separable.
Every normal projective $B$ satisfies $\mathcal H(B)$, because
every such $B'$ is projective.  Moreover $\mathcal H(B)$ implies
$\mathcal H(B_0)$ for every finite dominant normal $B_0\to B$,
by composition of finite morphisms.

\begin{proposition}[Proved: the exact base condition in the fibration argument]
\label{D:hereditary-base-generation}
In Proposition~\ref{D:projective-base-generation}, replace the
assumption that $B$ is projective by $\mathcal H(B)$.
Then the same conclusion holds, under either of its two
geometric generic-fibre hypotheses.
\end{proposition}
\begin{proof}
We specify the only point where the proof needs a property
of a projective base.  All arguments preceding the choice of
$Q'\subset Q$ use normality, properness, the fixed smooth
surjective base open, and the two valuation bounds.  They
therefore apply unchanged.  In particular the full wild
image at each of the finitely many base boundary divisors
is finite, including its inseparable residual contribution.

The normalization $B_1$ in the finite cover corresponding to
$Q'$ is finite and dominant over $B$, and normal and proper.
The same inertia intersection argument gives a surjection
\[
       Q(B_1,B_1\setminus B_1^\circ)\longrightarrow Q'.
\]
Property $\mathcal H(B)$ makes the source finitely generated.
The finite-extension argument for $Q$, and then the extension
by the finitely generated geometric image, proves the assertion.
No projectivity of $B_1$ is needed at this final step.
\end{proof}

\begin{theorem}[Proved: iterated fibrations with projective generic normalizations]
\label{D:iterated-fibration-generation}
Let $k$ be algebraically closed and suppose that
\[
 X_r\longrightarrow X_{r-1}\longrightarrow\cdots
       \longrightarrow X_0
\]
is a sequence of dominant morphisms of normal connected
proper $k$-schemes.  Suppose $X_0$ is projective and, for
each morphism in the sequence, the normalization of the
reduced geometric generic fibre is projective.  This
normalization may have more than one connected component.
Then $\mathcal H(X_i)$ holds for every $i$.

In particular the conclusion holds if every morphism has
generic fibre dimension at most one.  Thus for every finite
normal connected cover $Z\to X_r$ and every proper closed
subset $Y\subset Z$, the group $Q(Z,Y)$ is topologically
finitely generated.  If $Y$ has codimension at least two,
this is the full group $\pi_1(Z\setminus Y)$.
\end{theorem}
\begin{proof}
We prove the assertion by induction.  The projective case
gives $\mathcal H(X_0)$.  Suppose $\mathcal H(X_{i-1})$
has been proved, and fix a finite dominant morphism
$Z\to X_i$ with $Z$ normal and connected.  Stein factorization
of $Z\to X_{i-1}$ gives
\[
                 Z\longrightarrow B'\longrightarrow X_{i-1},
\]
where $B'$ is normal and the second morphism is finite.
The field $k(B')$ is algebraically closed in $k(Z)$.
Consequently the generic fibre is geometrically irreducible;
it need not be geometrically reduced.  This is the
relative-constant-field criterion already used in
Corollary~\ref{D:projective-image-corollaries}.

We verify that finite covers preserve the required generic
projectivity.  Over an algebraic closure of $k(B')$, each
normal component of the reduced generic fibre of $Z$
is finite over the normalization of a component of the
reduced geometric generic fibre of $X_i\to X_{i-1}$.
Indeed the original map is finite.  After reduction and
normalization the dominant map factors through the normal
target component; on affine rings it remains finite, since
an algebra finite over a smaller ring is finite over any
intermediate ring acting on it.  All these normalizations
are finite, as the schemes are of finite type over fields.
The target normal component is projective by hypothesis,
so the source normal component is projective as well.

Choose a finite purely inseparable extension $F_0/k(B')$
after which normalization makes the generic fibre
geometrically normal, using
\cite[Lemma 33.27.4, tag 0BXT]{C_StacksLogFibration}.
Normalize $B'$ in $F_0$ and $Z$ in $k(Z)F_0$, obtaining
$B_0$ and $Z_0$.  The first scheme is finite over $X_{i-1}$,
so $\mathcal H(B_0)$ follows from the induction hypothesis.
The second morphism $Z_0\to Z$ is finite and radicial:
a fixed $p$-power of every integral element of its function
field lies in $k(Z)$ and is integral over the normal local
ring downstairs, hence lies in that ring.  In characteristic
zero we take the trivial field extension.

The geometric generic fibre of $Z_0\to B_0$ is integral,
normal and projective by the preceding two paragraphs.
For $Y_0$ the inverse image of any proper closed $Y\subset Z$,
its fixed-boundary tame group is therefore finitely generated
by the projective case.  Proposition~\ref{D:hereditary-base-generation}
applied to $Z_0\to B_0$ gives finite generation of $Q(Z_0,Y_0)$.
Finite radicial base change preserves the category of finite
\'etale covers of the open and preserves tameness upon pullback
at the original boundary valuations, by
\cite[tags 04DY and 0EXY]{C_StacksLogFibration}.
It consequently gives a surjection
\[
                         Q(Z_0,Y_0)\longrightarrow Q(Z,Y).
\]
This proves $\mathcal H(X_i)$, including inseparable finite
covers, and completes the induction.

For the final sufficient condition, a normal proper scheme
over an algebraically closed field of dimension at most one
is projective, component by component, by
\cite[Lemma 33.43.4, section 0A22]{C_StacksLogFibration}.
Hence the required
generic projectivity is automatic for fibres of dimension
at most one.  If the boundary has codimension at least two,
there are no original divisorial wild inertia groups to kill,
so $Q(Z,Y)=\pi_1(Z\setminus Y)$.
\end{proof}

\begin{remark}
The sequence consists of morphisms defined on the entire
proper schemes.  Replacing a rational pencil by a morphism
on a modification does not give this hypothesis for the
original scheme: the additional horizontal exceptional
valuations can still have uncontrolled wild image.
\end{remark}
